% =============================================================================
% Chapter 7: Fit and Non-Concavity
% =============================================================================
%
% Version: 1.0 (January 2026)
%
% Purpose: Functional forms
%
% Primary Appendix: See appendix references below
% Secondary Appendices: G (Glossary), F (Examples)
%
% Prerequisites: Chapters 1-6
% Leads to: Chapter 8
%
% Chapter Type: B (Foundation)
% =============================================================================


%%%%%%%%%%%%%%%%%%%%%%%%%%%%%%%%%%%%%%%%%%%%%%%%%%%%%%%%%%%%%%%%%%%%%%%%%%%%%%%
% QUICK REFERENCE BOX
%%%%%%%%%%%%%%%%%%%%%%%%%%%%%%%%%%%%%%%%%%%%%%%%%%%%%%%%%%%%%%%%%%%%%%%%%%%%%%%
\begin{tcolorbox}[
    colback=gray!5!white,
    colframe=gray!75!black,
    title={\textbf{Begriffe in diesem Kapitel}},
    fonttitle=\bfseries
]
\small
\begin{tabular}{@{}ll@{}}
\textbf{Key Concept} & \textbf{Definition} \\
\midrule
Functional forms & See Section \ref{sec:ch7-intro} \\
\end{tabular}
\end{tcolorbox}



%%%%%%%%%%%%%%%%%%%%%%%%%%%%%%%%%%%%%%%%%%%%%%%%%%%%%%%%%%%%%%%%%%%%%%%%%%%%%%%
% INTUITION BOX
%%%%%%%%%%%%%%%%%%%%%%%%%%%%%%%%%%%%%%%%%%%%%%%%%%%%%%%%%%%%%%%%%%%%%%%%%%%%%%%
\begin{tcolorbox}[
    colback=blue!5!white,
    colframe=blue!75!black,
    title={\textbf{Intuition: A Motivating Example}},
    fonttitle=\bfseries
]
Consider \textbf{Anna}, a professional facing a decision that involves functional forms.

She initially evaluates the choice using standard criteria, but something feels incomplete.
\textbf{Thomas}, her colleague, suggests considering additional dimensions beyond the obvious.

This chapter explores what Anna discovers when she applies the EBF framework---how
functional forms interact with broader welfare considerations.

\medskip
\textit{By the end of this chapter, you will understand why Anna's initial analysis
was incomplete and how the EBF approach provides a more comprehensive view.}
\end{tcolorbox}



%%%%%%%%%%%%%%%%%%%%%%%%%%%%%%%%%%%%%%%%%%%%%%%%%%%%%%%%%%%%%%%%%%%%%%%%%%%%%%%
% CENTRAL QUESTION BOX
%%%%%%%%%%%%%%%%%%%%%%%%%%%%%%%%%%%%%%%%%%%%%%%%%%%%%%%%%%%%%%%%%%%%%%%%%%%%%%%
\begin{tcolorbox}[
    colback=red!5!white,
    colframe=red!75!black,
    title={\textbf{Central Question}},
    fonttitle=\bfseries
]
\Large
\centering
\textit{How does functional forms contribute to understanding human welfare and decision-making?}
\end{tcolorbox}



%%%%%%%%%%%%%%%%%%%%%%%%%%%%%%%%%%%%%%%%%%%%%%%%%%%%%%%%%%%%%%%%%%%%%%%%%%%%%%%
% CHAPTER OVERVIEW
%%%%%%%%%%%%%%%%%%%%%%%%%%%%%%%%%%%%%%%%%%%%%%%%%%%%%%%%%%%%%%%%%%%%%%%%%%%%%%%
\subsection*{Chapter Overview}

This chapter is organized as follows:

\begin{itemize}
\item \textbf{Section \ref{sec:ch7-intro}}: Introduction and motivation
\item \textbf{Section \ref{sec:ch7-theory}}: Theoretical foundations
\item \textbf{Section \ref{sec:ch7-applications}}: Applications and examples
\item \textbf{Section \ref{sec:ch7-summary}}: Summary and conclusions
\end{itemize}




%%%%%%%%%%%%%%%%%%%%%%%%%%%%%%%%%%%%%%%%%%%%%%%%%%%%%%%%%%%%%%%%%%%%%%%%%%%%%%%
% CHAPTER SCOPE BOX
%%%%%%%%%%%%%%%%%%%%%%%%%%%%%%%%%%%%%%%%%%%%%%%%%%%%%%%%%%%%%%%%%%%%%%%%%%%%%%%
\begin{tcolorbox}[
    colback=purple!5!white,
    colframe=purple!75!black,
    title={\textbf{Chapter Scope: Ziel / In-Scope / Out-of-Scope}},
    fonttitle=\bfseries
]

\textbf{Ziel (Objective):}
Develop the conceptual understanding of functional forms as a foundation for the EBF framework.

\vspace{0.3cm}
\textbf{In-Scope:}
\begin{itemize}[nosep]
    \item Conceptual introduction to functional forms
    \item Motivation and intuition
    \item Connection to subsequent chapters
    \item Key definitions and concepts
\end{itemize}

\vspace{0.3cm}
\textbf{Out-of-Scope (delegiert an Appendices):}
\begin{itemize}[nosep]
    \item Complete formal treatment $\rightarrow$ FORMAL appendices
    \item Empirical validation $\rightarrow$ METHOD appendices
    \item Domain applications $\rightarrow$ Application chapters
\end{itemize}

\end{tcolorbox}

%%%%%%%%%%%%%%%%%%%%%%%%%%%%%%%%%%%%%%%%%%%%%%%%%%%%%%%%%%%%%%%%%%%%%%%%%%%%%%%
% APPENDIX REFERENCES BOX
%%%%%%%%%%%%%%%%%%%%%%%%%%%%%%%%%%%%%%%%%%%%%%%%%%%%%%%%%%%%%%%%%%%%%%%%%%%%%%%
\begin{tcolorbox}[
    colback=yellow!10!white,
    colframe=orange!75!black,
    title={\textbf{Appendix References for This Chapter}},
    fonttitle=\bfseries
]
\small
\begin{tabular}{@{}lp{8cm}@{}}
\textbf{Appendix} & \textbf{Content} \\
\midrule
NC (FORMAL-NC) & Non-Concavity axioms NC-1 to NC-4, Proposition NC.1 \\
Y (DOMAIN-MR) & Milgrom-Roberts complementarity economics \\
G (REF-GLOSSARY) & Symbol definitions and terminology \\
\end{tabular}

\medskip
\textit{NC formalizes the $K$ vs.\ $Q$ distinction and the landscape of coherence.}
\end{tcolorbox}

\section{Complementarity, Fit, and Non-Concavity}
\label{sec:roberts}

Section~\ref{sec:cstar} established that $C^*$ is the structural attractor---
the configuration toward which systems gravitate.
But why do systems gravitate anywhere?
Why can't they wander freely across configuration space?

The answer lies in \emph{complementarity}: when elements reinforce each other,
payoff functions become non-concave, creating peaks and valleys.
Systems climb toward peaks and resist leaving them.

This section develops the connection to organizational economics (see Appendix~MIL for the Milgrom-Roberts supermodularity theory),
showing how \citet{roberts2004}'s concept of ``fit'' is precisely our coherence $K$,
and why the landscape metaphor is more than a metaphor.

\subsection{Roberts' Insight: Organizations as Systems of Complements}

\paragraph{The Problem Roberts Addressed.}
In the 1990s, organizational economics faced a puzzle:
why do successful firms look so different from each other,
yet each successful firm's elements fit together so well?

\begin{itemize}
  \item \textbf{Toyota:} Lean production, worker empowerment, continuous improvement, lifetime employment
  \item \textbf{Goldman Sachs:} Partnership culture, high risk tolerance, aggressive incentives, up-or-out
  \item \textbf{Southwest Airlines:} Point-to-point routes, single aircraft type, no frills, fun culture
\end{itemize}

Each is internally coherent. Each would fail if you mixed their elements.
Lean production with aggressive individual incentives? Disaster.
Partnership culture with standardized processes? Contradiction.

\paragraph{Roberts' Answer: Complementarity.}
\citet{roberts2004} explained this with \emph{complementarity}:
organizational elements are supermodular complements---
doing more of one increases the return to doing more of others.

\begin{equation}
  \frac{\partial^2 \Pi}{\partial x_i \partial x_j} > 0 
  \quad \text{for all } i \neq j
  \label{eq:supermodularity}
\end{equation}

where $x_i$ are organizational choices: strategy, structure, processes, 
people, culture, incentives.

\paragraph{Example: Strategy and Structure.}
Consider two strategic choices:
\begin{itemize}
  \item \textbf{Innovation strategy:} Develop new products, enter new markets
  \item \textbf{Efficiency strategy:} Cut costs, standardize, optimize
\end{itemize}

And two structural choices:
\begin{itemize}
  \item \textbf{Flat structure:} Few hierarchy levels, autonomous teams
  \item \textbf{Hierarchical structure:} Many levels, clear command chains
\end{itemize}

The cross-derivatives are positive:
\begin{align}
  \frac{\partial^2 \Pi}{\partial \text{Innovation} \, \partial \text{Flat}} &> 0 
  \quad \text{(Innovation needs autonomy)} \\
  \frac{\partial^2 \Pi}{\partial \text{Efficiency} \, \partial \text{Hierarchy}} &> 0 
  \quad \text{(Efficiency needs control)}
\end{align}

But the cross-configurations are penalized:
\begin{align}
  \frac{\partial^2 \Pi}{\partial \text{Innovation} \, \partial \text{Hierarchy}} &< 0 
  \quad \text{(Innovation stifled by bureaucracy)} \\
  \frac{\partial^2 \Pi}{\partial \text{Efficiency} \, \partial \text{Flat}} &< 0 
  \quad \text{(Efficiency lost to coordination problems)}
\end{align}

\subsection{Non-Concavity: Why Complementarity Creates Peaks and Valleys}

\paragraph{The Mathematical Structure.}
Under complementarity, the payoff function $\Pi(x_1, \ldots, x_n)$ is 
\emph{supermodular}: its Hessian has positive off-diagonal elements.

This has a crucial implication: the function is \emph{not globally concave}.

\begin{proposition}[Non-Concavity from Complementarity]
\label{prop:nonconcavity}
If $\partial^2 \Pi / \partial x_i \partial x_j > 0$ for all $i \neq j$,
then $\Pi$ is not globally concave.
The payoff landscape has multiple local maxima.
\end{proposition}

\begin{proof}
The Hessian matrix $H$ of $\Pi$ has:
\begin{itemize}
  \item Diagonal elements $H_{ii} = \partial^2 \Pi / \partial x_i^2 < 0$ (concavity in own choice)
  \item Off-diagonal elements $H_{ij} = \partial^2 \Pi / \partial x_i \partial x_j > 0$ (complementarity)
\end{itemize}

For global concavity, $H$ must be negative semi-definite: all eigenvalues $\leq 0$.
But with positive off-diagonal elements, the matrix has at least one positive eigenvalue.
Hence $\Pi$ is not concave.
\end{proof}

\paragraph{The Landscape Metaphor.}
Think of $\Pi$ as a landscape over configuration space:
\begin{itemize}
  \item \textbf{Peaks:} Coherent configurations where all elements fit together
  \item \textbf{Valleys:} Incoherent configurations where elements clash
  \item \textbf{Ridges:} Partial coherence---some elements fit, others don't
\end{itemize}

\paragraph{Example: The Two-Peak Landscape.}
With two choices (strategy $s$, structure $\sigma$), each in $[0,1]$:

\begin{center}
\begin{tabular}{ccc}
& $\sigma = 0$ (Hierarchy) & $\sigma = 1$ (Flat) \\
\hline
$s = 0$ (Efficiency) & \textbf{Peak A}: $\Pi = \beta$ & Valley: $\Pi < 0$ \\
$s = 1$ (Innovation) & Valley: $\Pi < 0$ & \textbf{Peak B}: $\Pi = \alpha$ \\
\end{tabular}
\end{center}

Two peaks (A and B), with a valley between them.
An organization at (0.5, 0.5)---half innovative, half efficient---is in the valley.

\subsection{Fit as Coherence}

\paragraph{The Roberts-Framework Connection.}
\citet{roberts2004} defines \textbf{fit} as the degree to which
organizational elements reinforce each other.
This maps directly to our framework:

\begin{center}
\renewcommand{\arraystretch}{1.3}
\begin{tabular}{lll}
\hline
\textbf{Roberts (2004)} & \textbf{Present Framework} & \textbf{Meaning} \\
\hline
Fit & $K \approx 1$ & At a local maximum \\
Misfit & $K < 1$ & In the valley between peaks \\
Configuration quality & $Q(C^*)$ & Height of the peak \\
Multiple configurations & Multiple $C^*$ & Multiple local maxima \\
Strategy-structure alignment & $s \approx \sigma$ & On the diagonal \\
\hline
\end{tabular}
\end{center}

\paragraph{Example: Southwest Airlines.}
Southwest's configuration:
\begin{itemize}
  \item Strategy: Low cost, high frequency
  \item Structure: Decentralized, empowered employees
  \item Processes: Fast turnaround, no assigned seats
  \item People: Hired for attitude, trained for skill
  \item Culture: Fun, team-oriented
  \item Incentives: Profit sharing, no layoffs
\end{itemize}

Every element reinforces every other. Change one, and the system breaks.
This is $K \approx 1$: high coherence.

\paragraph{Example: Failed Transformation at Continental.}
In the 1990s, Continental Airlines tried to copy Southwest:
\begin{itemize}
  \item Adopted low-cost strategy (changed $s$)
  \item Kept hierarchical structure (didn't change $\sigma$)
  \item Kept old processes, people, culture, incentives
\end{itemize}

Result: Chaos. The strategy required different structure, but structure didn't change.
$K \ll 1$: Low coherence. Employees didn't know whether to innovate or follow rules.
Continental retreated and eventually found success with a \emph{different} coherent configuration.

\subsection{Why Transitions Are So Difficult}

\paragraph{The Valley Problem.}
To move from Peak A to Peak B, an organization must:
\begin{enumerate}
  \item Leave Peak A (coherence falls)
  \item Cross the valley (incoherence, pain)
  \item Climb to Peak B (coherence restored)
\end{enumerate}

\begin{equation}
  \text{Transition path: } 
  C^*_A \xrightarrow{K \downarrow} \text{Valley} 
        \xrightarrow{K \uparrow} C^*_B
\end{equation}

\paragraph{Why Partial Change Fails.}
The complementarity structure means you can't change elements one at a time.
If you change strategy but not structure, you move \emph{sideways}---
not toward the new peak, but into the valley.

\begin{center}
\renewcommand{\arraystretch}{1.2}
\begin{tabular}{lcccc}
\hline
\textbf{Change} & $\Delta s$ & $\Delta \sigma$ & Direction & Result \\
\hline
Full transformation & $+0.8$ & $+0.8$ & Diagonal & Reach Peak B \\
Strategy only & $+0.8$ & $0$ & Horizontal & Stuck in valley \\
Structure only & $0$ & $+0.8$ & Vertical & Stuck in valley \\
Gradual (sequential) & $+0.8$, then $+0.8$ & & Two valleys & Worse than staying \\
\hline
\end{tabular}
\end{center}

\paragraph{Example: Digital Transformation.}
A traditional manufacturing firm wants to become a digital innovator.

\textbf{The wrong way (sequential):}
\begin{enumerate}
  \item Year 1: Announce digital strategy
  \item Year 2: Start changing processes
  \item Year 3: Realize culture doesn't fit
  \item Year 4: Hire new people
  \item Year 5: Change incentives
  \item Years 6-10: Still in the valley, key people have left
\end{enumerate}

\textbf{The right way (simultaneous):}
\begin{enumerate}
  \item Year 0: Plan complete transformation
  \item Year 1: Change strategy, structure, processes, hiring, culture, incentives \emph{together}
  \item Years 2-3: Accept low $K$, support the transition
  \item Year 4: New coherent configuration emerges
\end{enumerate}

The second path is harder in the short run but reaches the new peak.

\subsection{Multiple Peaks: Which Configuration Is Best?}

\paragraph{K vs. Q Again.}
Both Peak A (Efficiency) and Peak B (Innovation) have $K \approx 1$.
Both are coherent. But they differ in $Q$---the height of the peak.

Which peak is higher depends on context $\Psi$:
\begin{itemize}
  \item Stable industry, commodity product: Peak A (Efficiency) is higher
  \item Turbulent industry, differentiated product: Peak B (Innovation) is higher
\end{itemize}

\paragraph{Example: Kodak's Coherent Decline.}
Kodak in 2000 was highly coherent ($K \approx 1$):
\begin{itemize}
  \item Strategy: Dominate film photography
  \item Structure: Vertically integrated
  \item Processes: Optimized for chemical film
  \item Culture: Engineering excellence in film
\end{itemize}

But context changed: digital photography emerged.
Kodak's peak (film excellence) became a valley in the new landscape.
High $K$, but $Q$ collapsed because $\Psi$ shifted.

\paragraph{The Trap of Coherent Obsolescence.}
A firm can be perfectly coherent and still fail---
if it's coherent in a configuration that the market no longer rewards.

This is the $K$ vs. $Q$ distinction in action:
\begin{itemize}
  \item $K$ measures fit with self: Are our elements aligned?
  \item $Q$ measures fit with world: Is our configuration valued?
\end{itemize}

\subsection{From Organizations to All Levels}

\paragraph{Roberts' Insight Generalizes.}
The complementarity-fit-non-concavity logic applies beyond firms:

\begin{center}
\renewcommand{\arraystretch}{1.3}
\begin{tabular}{llll}
\hline
\textbf{Level} & \textbf{Elements} & \textbf{Fit Means} & \textbf{Example} \\
\hline
Individual & Values, goals, behavior, habits & Integrity & Live your values \\
Household & Roles, expectations, resources & Family harmony & Shared parenting \\
Organization & Strategy, structure, culture & Competitive fit & Southwest \\
Region & Economy, education, infrastructure & Economic coherence & Silicon Valley \\
Nation & Institutions, norms, policies & Institutional fit & Nordic model \\
Global & Trade, climate, security regimes & International order & Bretton Woods \\
\hline
\end{tabular}
\end{center}

At every level, complementarity creates peaks and valleys.
At every level, partial change lands you in the valley.
At every level, $K$ and $Q$ are distinct.

\subsection{Implications for Strategy}

\paragraph{Diagnostic Questions.}
Before any change, ask:
\begin{enumerate}
  \item Where are we? (Which peak, or are we in a valley?)
  \item Where do we want to be? (Which peak is best given $\Psi$?)
  \item How deep is the valley? (Transition cost)
  \item Can we change everything together? (Transformation capacity)
\end{enumerate}

\paragraph{Strategic Options.}
\begin{itemize}
  \item \textbf{Optimize:} Stay on current peak, improve $Q$ marginally
  \item \textbf{Transform:} Full-stack change to different peak
  \item \textbf{Hybrid:} Rarely works---usually lands in valley
  \item \textbf{Exit:} If valley too deep and transformation impossible
\end{itemize}

\paragraph{The Key Insight.}
Complementarity is both blessing and curse:
\begin{itemize}
  \item Blessing: Coherent configurations are stable (hard to dislodge)
  \item Curse: Coherent configurations are traps (hard to leave)
\end{itemize}

Understanding this structure is the first step toward navigating it.

\subsection{Why Navigation Requires Calibration, Not Description}

\paragraph{The Strategic Question LLMs Cannot Answer.}
A firm facing digital disruption asks: ``Should we transform?''
This requires answering:
\begin{enumerate}
    \item How deep is the valley between our current peak and the digital peak?
    \item What is the height of the digital peak in our specific context?
    \item Do we have the resources to sustain the transition?
\end{enumerate}

These are \emph{quantitative} questions about a specific landscape.
An LLM can explain the concepts---peaks, valleys, complementarity, fit---fluently.
But it cannot navigate the landscape because it doesn't know its shape.

\paragraph{Why Text Cannot Map the Landscape.}
The payoff landscape $\Pi(x_1, \ldots, x_n; \Psi)$ depends on:
\begin{itemize}
    \item The complementarity structure $C_{ij}$: How strongly do elements reinforce each other?
    \item The context $\Psi$: What does the market reward?
    \item The transition dynamics: How fast can capabilities be rebuilt?
\end{itemize}

None of these can be extracted from text about organizational transformation.
The literature on digital transformation describes \emph{that} firms struggle,
documents \emph{cases} of success and failure, and proposes \emph{principles}
for managing change. But it does not quantify:
\begin{itemize}
    \item The valley depth for a manufacturing firm versus a service firm
    \item How context $\Psi$ shifts the relative heights of peaks
    \item The interaction between transformation speed and survival probability
\end{itemize}

\paragraph{Example: Kodak's Optimal Response.}
The Kodak case is often told as a cautionary tale: ``They should have transformed.''
But what was the actual landscape?

\begin{center}
\small
\renewcommand{\arraystretch}{1.2}
\begin{tabular}{lcc}
\toprule
\textbf{Parameter} & \textbf{Qualitative} & \textbf{Quantitative (Calibrated)} \\
\midrule
Film peak height (2000) & ``Declining'' & $Q_{film} \approx 0.7$ (still profitable) \\
Digital peak height (2000) & ``Rising'' & $Q_{digital} \approx 0.3$ (negative margins) \\
Valley depth & ``Deep'' & $\Delta \Pi \approx -\$5B$ over 5 years \\
Transition probability & ``Uncertain'' & $P(success) \approx 0.2$ \\
\bottomrule
\end{tabular}
\end{center}

With calibrated parameters, the optimal strategy might have been:
\emph{harvest the film business while funding a separate digital entity}---
not the all-in transformation that conventional wisdom suggests.
The point is not that this was correct, but that the \emph{answer depends on numbers}
that text-based analysis cannot provide.

\paragraph{The Calibration Imperative for Strategy.}
Section~4.X established that behavioral prediction requires calibration on data.
The same logic applies to strategic navigation:
\begin{itemize}
    \item Predicting behavior: Calibrate $C^*(\Psi)$ on experiments
    \item Navigating landscapes: Calibrate $\Pi(x; \Psi)$ on firm performance data
\end{itemize}

The complementarity-fit-non-concavity framework provides the \emph{structure}:
we know landscapes have peaks and valleys, transitions are costly, and fit matters.
But the \emph{parameters}---how deep are valleys, how high are peaks---must be
estimated from data, not derived from principles or learned from text.

This is why strategic consulting requires industry-specific analysis,
why transformation success rates vary by sector, and why generic advice
(``embrace digital'') fails. The landscape is context-dependent,
and context must be measured, not assumed.

%------------------------------------------------------------------------------
% SUMMARY
%------------------------------------------------------------------------------
\section{Summary}
\label{sec:ch7-summary}

\begin{tcolorbox}[
    colback=blue!5!white,
    colframe=blue!75!black,
    title={\textbf{Chapter 7 Summary: Complementarity, Fit, and Non-Concavity}}
]

\textbf{Core Insight:} Complementarity ($\gamma > 0$) creates non-concave payoff landscapes with peaks ($K \approx 1$) and valleys ($K \ll 1$), making transitions costly and stability path-dependent.

\textbf{Key Results:}
\begin{enumerate}[nosep]
    \item \textbf{Non-Concavity:} When $\gamma_{ij} > 0$, the payoff function $\Pi(C)$ is non-concave
    \item \textbf{Multiple Equilibria:} Systems can be stable at different $K$ levels---not all equilibria are equal
    \item \textbf{Transition Costs:} Moving from low-$K$ to high-$K$ requires crossing valleys
    \item \textbf{Fit:} Coherent configurations ($K \approx 1$) are locally stable but require coordination
\end{enumerate}

\textbf{Transition:} Chapter 8 provides the mathematical formalization of these concepts.
\end{tcolorbox}

\subsection{What Comes Next}
\label{sec:ch7-next}

With the non-concavity principle established, subsequent chapters formalize and apply it:

\begin{enumerate}[nosep]
    \item \textbf{Chapter 8:} Mathematical foundations---formal treatment of dynamics
    \item \textbf{Chapter 9:} Context as endogenous---how $\Psi$ modulates the landscape
    \item \textbf{Chapter 10:} Die Nutzenarchitektur---applying fit to welfare analysis
\end{enumerate}

%%%%%%%%%%%%%%%%%%%%%%%%%%%%%%%%%%%%%%%%%%%%%%%%%%%%%%%%%%%%%%%%%%%%%%%%%%%%%%%
% READING PATH BOX
%%%%%%%%%%%%%%%%%%%%%%%%%%%%%%%%%%%%%%%%%%%%%%%%%%%%%%%%%%%%%%%%%%%%%%%%%%%%%%%
\begin{tcolorbox}[
    colback=green!5!white,
    colframe=green!75!black,
    title={\textbf{Reading Path}},
    fonttitle=\bfseries
]
\begin{tabular}{@{}ll@{}}
\textbf{Previous:} & Chapter 6 \\
\textbf{Next:} & Chapter 8 \\
\textbf{Deep Dive:} & See Appendix G (Glossary) for definitions \\
\end{tabular}

\medskip
\textit{For readers short on time: Focus on the Intuition Box and Summary section.}
\end{tcolorbox}

